\chapter{Phụ lục}

\section{Phân chia công việc}

Phần này trình bày chi tiết về phân chia công việc giữa các thành viên trong nhóm cho giai đoạn 1 của dự án (từ Chương 1 đến Chương 6). Công việc được phân chia một cách công bằng dựa trên khả năng và sở trường của từng thành viên, đồng thời đảm bảo tính hợp tác và hỗ trợ lẫn nhau trong quá trình phát triển dự án.

\begin{longtable}{|p{1cm}|p{3.5cm}|p{10cm}|}
\hline
\textbf{STT} & \textbf{Họ và tên} & \textbf{Công việc} \\
\hline
\endfirsthead

\hline
\textbf{STT} & \textbf{Họ và tên} & \textbf{Công việc} \\
\hline
\endhead

\hline
\endfoot

\hline
\endlastfoot

1 & Trần Hoàng Sơn & 1.1. Viết báo cáo Chương 1: Giới thiệu (động lực phát triển, phạm vi và mục tiêu dự án, cấu trúc báo cáo) \\
& & 2.1.1. Nghiên cứu và phân tích các nền tảng social commerce hiện có (TikTok Shop, Facebook Marketplace) \\
& & 2.1.2. Phân tích các hệ thống quản lý thương mại điện tử để tham khảo \\
& & 3.1.3. Thiết kế Use Case Diagram cho Module 3: Hoạt động Thương mại (Products, Cart, Orders, Reviews) \\
& & 3.2.3. Vẽ Activity Diagram cho các luồng E-commerce (thêm vào giỏ hàng, thanh toán, đánh giá sản phẩm) \\
& & 3.3.3. Thiết kế Sequence Diagram cho các chức năng thương mại (đặt hàng, quản lý đơn hàng) \\
& & 3.4.3. Thiết kế schema database cho các bảng liên quan đến E-commerce (Products, Cart, Orders, Reviews) \\
& & 4.3. Nghiên cứu và so sánh các hệ quản trị cơ sở dữ liệu (PostgreSQL, MySQL, MongoDB) cho Chương 4 \\
& & 5.1.3. Thiết kế giao diện UI/UX trên Figma cho Module 3: Thương mại (trang sản phẩm, giỏ hàng, thanh toán, quản lý đơn hàng) \\
& & 5.2.3. Viết báo cáo Chương 5: Triển khai Hệ thống (phần thiết kế UI/UX Module 3) \\
& & 6.1. Đóng góp vào viết báo cáo Chương 6: Tổng kết và Kế hoạch Tương lai \\
& & 6.2. Review và chỉnh sửa toàn bộ tài liệu báo cáo \\
\hline

2 & Phạm Lê Thanh & 2.1.1. Nghiên cứu và phân tích các nền tảng social commerce hiện có (Instagram Shopping, Pinterest Shopping) \\
& & 2.2. Viết báo cáo Chương 2: Thu thập và Phân tích Yêu cầu (phân tích các hệ thống liên quan, xác định actors và functional requirements) \\
& & 3.1.2. Thiết kế Use Case Diagram cho Module 2: Tương tác Xã hội (Posts, Comments, Follow, Messages, Groups) \\
& & 3.2.2. Vẽ Activity Diagram cho các luồng Social (tạo bài đăng, tương tác, nhắn tin, tham gia nhóm) \\
& & 3.3.2. Thiết kế Sequence Diagram cho các chức năng xã hội (đăng bài, bình luận, theo dõi, chat real-time) \\
& & 3.4.2. Thiết kế schema database cho các bảng liên quan đến Social Features (Posts, Comments, Messages, Notifications) \\
& & 4.1. Nghiên cứu và so sánh các công nghệ Frontend (React, Vue, Angular) cho Chương 4 \\
& & 4.4. Nghiên cứu và so sánh các công nghệ Real-time (Socket.io, WebSocket, Server-Sent Events) cho Chương 4 \\
& & 4.5. Viết báo cáo Chương 4: Nghiên cứu Công nghệ (phần Frontend và Real-time) \\
& & 5.1.2. Thiết kế giao diện UI/UX trên Figma cho Module 2: Tương tác Xã hội (feed, tạo bài đăng, nhắn tin, nhóm) \\
& & 5.1.4. Thiết kế giao diện UI/UX cho Module 6, 7: Lên lịch đăng bài và Hệ thống Thông báo \\
& & 6.1. Đóng góp vào viết báo cáo Chương 6: Tổng kết và Kế hoạch Tương lai \\
\hline

3 & Nguyễn Nhật Tân & 2.1.1. Nghiên cứu và phân tích các nền tảng social commerce hiện có (Shopee, Lazada) \\
& & 2.2.1. Xác định và mô tả chi tiết các Actors (Visitor, User, Buyer, Seller, Admin) cho Chương 2 \\
& & 2.3. Đặc tả các yêu cầu phi chức năng (Performance, Security, Availability, Usability, Scalability, Real-time) cho Chương 2 \\
& & 3.1.1. Thiết kế Use Case Diagram cho Module 1: Quản lý Người dùng và Module 4: Hệ thống Quản trị \\
& & 3.2.1. Vẽ Activity Diagram cho các luồng Authentication (đăng ký, đăng nhập, 2FA, quên mật khẩu) \\
& & 3.3.1. Thiết kế Sequence Diagram cho các chức năng quản trị (quản lý người dùng, kiểm duyệt nội dung, xử lý báo cáo) \\
& & 3.4.1. Thiết kế ERD tổng thể và schema database cho các bảng liên quan đến User Management và Admin \\
& & 3.5. Thiết kế Architecture Diagram tổng thể của hệ thống \\
& & 3.6. Thiết kế Component Diagram và Deployment Diagram \\
& & 3.7. Viết báo cáo Chương 3: Thiết kế Hệ thống (phần Architecture Diagram, Component Diagram, Deployment Diagram) \\
& & 4.2. Nghiên cứu và so sánh các công nghệ Backend (Node.js/Express, Django, Spring Boot) cho Chương 4 \\
& & 4.6. Nghiên cứu và so sánh các công nghệ AI (Google Gemini, OpenAI GPT, Claude) cho Chương 4 \\
& & 5.1.1. Thiết kế giao diện UI/UX trên Figma cho Module 1: Quản lý Người dùng (đăng ký, đăng nhập, profile, settings) \\
& & 5.1.5. Thiết kế giao diện UI/UX trên Figma cho Module 4: Hệ thống Quản trị (dashboard, quản lý người dùng, báo cáo) \\
& & 6.1. Đóng góp vào viết báo cáo Chương 6: Tổng kết và Kế hoạch Tương lai \\
\hline

\end{longtable}

\section{Ghi chú về phân chia công việc}

\begin{itemize}
    \item Công việc được phân chia dựa trên sở trường và khả năng của từng thành viên, đảm bảo tính công bằng và hiệu quả.
    \item Các công việc quan trọng như thiết kế ERD, Architecture Diagram được thực hiện chung để đảm bảo tính nhất quán và chất lượng.
    \item Việc phân chia công việc được điều chỉnh linh hoạt trong quá trình phát triển dự án để đảm bảo tiến độ và chất lượng.
    \item Tất cả các thành viên đều tham gia vào quá trình review và chỉnh sửa để đảm bảo chất lượng tài liệu.
    \item Các module phức tạp được phân chia để tận dụng điểm mạnh của từng thành viên, đồng thời đảm bảo tính hợp tác.
\end{itemize}

